%%%%%%%%%%%%%%%%%%%%%%%%%%%%%%%%%%%%%%%%%
% Professional Formal Letter
% LaTeX Template
% Version 2.0 (12/2/17)
%
% This template originates from:
% http://www.LaTeXTemplates.com
%
% Authors:
% Brian Moses
% Vel (vel@LaTeXTemplates.com)
%
% License:
% CC BY-NC-SA 3.0 (http://creativecommons.org/licenses/by-nc-sa/3.0/)
%
%%%%%%%%%%%%%%%%%%%%%%%%%%%%%%%%%%%%%%%%%

%----------------------------------------------------------------------------------------
%	PACKAGES AND OTHER DOCUMENT CONFIGURATIONS
%----------------------------------------------------------------------------------------

\documentclass[11pt, a4paper]{letter} % Set the font size (10pt, 11pt and 12pt) and paper size (letterpaper, a4paper, etc)

%%%%%%%%%%%%%%%%%%%%%%%%%%%%%%%%%%%%%%%%%
% Professional Formal Letter
% Structure Specification File
% Version 1.0 (12/2/17)
%
% This file originates from:
% http://www.LaTeXTemplates.com
%
% Authors:
% Brian Moses
% Vel (vel@LaTeXTemplates.com)
%
% License:
% CC BY-NC-SA 3.0 (http://creativecommons.org/licenses/by-nc-sa/3.0/)
%
%%%%%%%%%%%%%%%%%%%%%%%%%%%%%%%%%%%%%%%%%

%----------------------------------------------------------------------------------------
%	PACKAGES AND OTHER DOCUMENT CONFIGURATIONS
%----------------------------------------------------------------------------------------

\usepackage{graphicx} % Required for including pictures

\usepackage[T1]{fontenc} % Output font encoding for international characters
\usepackage[utf8]{inputenc} % Required for inputting international characters
\usepackage{gfsdidot} % Use the GFS Didot font: http://www.tug.dk/FontCatalogue/gfsdidot/
\usepackage{microtype} % Improves typography

\pagestyle{empty} % Suppress headers and footers

\setlength\parindent{1cm} % Paragraph indentation

% Create a new command for the horizontal rule in the document which allows thickness specification
\makeatletter
\newcommand{\vhrulefill}[1]{\leavevmode\leaders\hrule\@height#1\hfill \kern\z@}
\makeatother

%----------------------------------------------------------------------------------------
%	DOCUMENT MARGINS
%----------------------------------------------------------------------------------------

\usepackage{geometry} % Required for adjusting page dimensions

\geometry{
	top=1cm, % Top margin
	bottom=1.5cm, % Bottom margin
	left=3cm, % Left margin
	right=3cm, % Right margin
	%showframe, % Uncomment to show how the type block is set on the page
}

%----------------------------------------------------------------------------------------
%	DEFINE CUSTOM COMMANDS
%----------------------------------------------------------------------------------------

\newcommand{\logo}[1]{\renewcommand{\logo}{#1}}

\newcommand{\Who}[1]{\renewcommand{\Who}{#1}}
\newcommand{\Title}[1]{\renewcommand{\Title}{#1}}

\newcommand{\headerlineone}[1]{\renewcommand{\headerlineone}{#1}}
\newcommand{\headerlinetwo}[1]{\renewcommand{\headerlinetwo}{#1}}

\newcommand{\authordetails}[1]{\renewcommand{\authordetails}{#1}}

%----------------------------------------------------------------------------------------
%	AUTHOR DETAILS STRUCTURE
%----------------------------------------------------------------------------------------

\newcommand{\authordetailsblock}{
	\hspace{\fill} % Move the author details to the far right
	\parbox[t]{0.48\textwidth}{ % Box holding the author details; width value specifies where it starts and ends, increase to move details left
		\footnotesize % Use a smaller font size for the details
		\Who\\ % Author name
		\textit{\authordetails} % The author details text, all italicised
	}
}

%----------------------------------------------------------------------------------------
%	HEADER STRUCTURE
%----------------------------------------------------------------------------------------

\address{
	\includegraphics[width=2.5in]{\logo} % Include the logo of author institution
	\hspace{0.82\textwidth} % Position of the institution logo, increase to move left, decrease to move right
	\vskip -0.05\textheight~\\ % Position of the large header text in relation to the institution logo, increase to move down, decrease to move up
	\Large\hspace{0.2\textwidth}\headerlineone\hfill ~\\[0.006\textheight] % First line of institution name, adjust hspace if your logo is wide
	\hspace{0.2\textwidth}\headerlinetwo\hfill \normalsize % Second line of institution name, adjust hspace if your logo is wide
% 	\makebox[0ex][r]{\textbf{\Who\Title}}\hspace{0.01\textwidth} % Print author name and title with a little whitespace to the right
	~\\[-0.01\textheight] % Reduce the whitespace above the horizontal rule
	\hspace{0.2\textwidth}\vhrulefill{1pt} \\ % Horizontal rule, adjust hspace if your logo is wide and \vhrulefill for the thickness of the rule
	\authordetailsblock % Include the letter author's details on the right side of the page under the horizontal rule
	\hspace{-0.25\textwidth} % Horizontal position of the author details block, increase to move left, decrease to move right
	\vspace{-0.1\textheight} % Move the date and letter content up for a more compact look
}

%----------------------------------------------------------------------------------------
%	COMPOSE THE ENTIRE HEADER
%----------------------------------------------------------------------------------------

\renewcommand{\opening}[1]{
	{\centering\fromaddress\vspace{0.05\textheight} \\ % Print the header and from address here, add whitespace to move date down
	\hspace*{\longindentation}\today\hspace*{\fill}\par} % Print today's date, remove \today to not display it
	{\raggedright \toname \\ \toaddress \par} % Print the to name and address
	\vspace{1cm} % White space after the to address
	\noindent #1 % Print the opening line
}

%----------------------------------------------------------------------------------------
%	SIGNATURE STRUCTURE
%----------------------------------------------------------------------------------------

\signature{\Who\Title} % The signature is a combination of the author's name and title

\renewcommand{\closing}[1]{
	\vspace{2.5mm} % Some whitespace after the letter content and before the signature
	\noindent % Stop paragraph indentation
	\hspace*{\longindentation} % Move the signature right to the value of \longindentation
	\parbox{\indentedwidth}{
		\raggedright
		#1 % Print the signature text
		\vskip 1.65cm % Whitespace between the closing text and author's name for a physical signature
		\fromsig \\% Prints the value of \signature{}, i.e. author name and title
		\includegraphics[scale=0.25]{lmfcarvalho_sign.jpg}
	}
}
 % Include the file that specifies the document structure

%\longindentation=0pt % Un-commenting this line will push the closing "Sincerely," and date to the left of the page

%----------------------------------------------------------------------------------------
%	YOUR INFORMATION
%----------------------------------------------------------------------------------------


\Who{Luiz Max de Carvalho} % Your name

\Title{, PhD} % Your title, leave blank for no title
\authordetails{
	School of Applied Mathematics\\ % Your department/institution
	Praia de Botafogo, 190\\ % Your address
	Rio de Janeiro, RJ, 22250-900\\ % Your city, zip code, country, etc
	Email: lmax.fgv@gmail.com \\ % Your email address
	Phone: +55 21 3799-2348 \\ % Your phone number
% 	URL: LaTeXTemplates.com % Your URL
}

%----------------------------------------------------------------------------------------
%	HEADER CONTENTS
%----------------------------------------------------------------------------------------
\logo{emap.png}
% \logo{Marca_FGV_EMAp_colorida.png}

\headerlinetwo{Getúlio Vargas Foundation (FGV)} % Top header line, leave blank if you only want the bottom line

% \headerlinetwo{School of Applied Mathematics (EMAp)} % Bottom header line

%----------------------------------------------------------------------------------------

\begin{document}

%----------------------------------------------------------------------------------------
%	TO ADDRESS
%----------------------------------------------------------------------------------------

\begin{letter}{
	Prof. Jane-Ling Wang\\
	Department of Statistics\\
	University of California\\
	Davis, CA 95616, USA
}

%----------------------------------------------------------------------------------------
%	LETTER CONTENT
%----------------------------------------------------------------------------------------

\opening{Dear Professor Wang,}

I would like to submit the attached manuscript for consideration for publication in~\textit{JASA}.
The paper is being submitted for the first time.
It extends the approach of Poole \& Raftery (2000)\footnote{Poole D, Raftery AE. Inference for deterministic simulation models: the Bayesian melding approach. Journal of the American Statistical Association. 2000 Dec 1;95(452):1244-55.} to accommodate varying weights in a logarithmic pool (LP).
Instead of focusing on Bayesian melding only, however, we offer a more general study of log-linear mixtures and their statistical applications.
A number of caveats are also transparently discussed.

In the same fashion as other~\textit{JASA} papers such as~Gneiting \& Raftery (2007)\footnote{Gneiting T, Raftery AE. Strictly proper scoring rules, prediction, and estimation. Journal of the American statistical Association. 2007 Mar 1;102(477):359-78.}, we offer a mix of review of important results and new ones.
We offer a concise but inclusive exposition of the main properties of logarithmic pooling, as well as some new results such as log-pooling for exponential families, unique log-concavity of LP and non-uniqueness of the entropy-optimised weights.
We then bring the main innovation of the paper, which are the hierarchical priors on the weights, discuss prior choice in this setting and explore issues of identifiability and interpretation in detail.
To complete exposition, we provide a wealth of applications, from simple probability estimation to meta-analysis to Bayesian melding for deterministic models -- including epidemic models.
The paper also leaves open many new avenues of research, such as using logarithmic pooling for Bayesian model averaging and eliciting priors for the weights from external information.

This would be of interest to the readership of~\textit{JASA} because it sheds light on an overlooked part of the statistical literature while offering motivation for the use of logarithmic pooling (i.e, its mathematical properties) and  a clear exposition of the potential applications in a broad class of problems.
The paper is the culmination of over five years of work.
We are indebted to the people listed in the Acknowledgments section of the paper, who helped us solidify ideas and hone many of the arguments/experiments in the paper over the years. 
This is some of my best work and I would be glad if you would consider it for publication.

\closing{Cordially,}

%----------------------------------------------------------------------------------------
%	OPTIONAL FOOTNOTE
%----------------------------------------------------------------------------------------

% Uncomment the 4 lines below to print a footnote with custom text
%\def\thefootnote{}
%\def\footnoterule{\hrule}
%\footnotetext{\hspace*{\fill}{\footnotesize\em Footnote text}}
%\def\thefootnote{\arabic{footnote}}

%----------------------------------------------------------------------------------------

\end{letter}

\end{document}
